\documentclass[11pt]{amsart} 
\usepackage{amsmath,amssymb,a4wide} 
\usepackage{color}
\usepackage{mathrsfs}
\usepackage{dsfont}
% \usepackage{showkeys}
\usepackage{hyperref}
\usepackage{mathabx}
\usepackage{tikz,pgfplots}
\newtheorem{theorem}{Theorem} 
\newtheorem{lemma}[theorem]{Lemma} 
\newtheorem{prop}[theorem]{Proposition} 
\newtheorem{remark}[theorem]{Remark} 
\newtheorem{corollary}[theorem]{Corollary}
\newtheorem{defi}[theorem]{Definition}

\newcommand{\pr}{\noindent  {\bf \textit{Proof : }}}

\newcommand{\cqfd}{{\nobreak\hfil\penalty50\hskip2em\hbox{}\nobreak\hfil
$\square$\qquad\parfillskip=0pt\finalhyphendemerits=0\par\medskip}}

\newcommand{\red}{\textcolor{red}}

\newcommand{\dis}{\displaystyle}
\newcommand{\dt}{\delta t}

\newcommand{\tore}{\Omega}

\newcommand{\Z}{{\mathbb Z}} 
\newcommand{\R}{{\mathbb R}} 
\newcommand{\C}{{\mathbb C}} 
\newcommand{\N}{{\mathbb N}} 
\newcommand{\TD}{{\mathbb T}} 
\newcommand{\T}{{\mathbb T}} 
\newcommand{\E}{{\mathcal E}}
\newcommand{\dd}{{\rm d}}
\newcommand{\ii}{{\rm i}}
\newcommand{\name}{{\rm MM\_Propagation }}

\let\pa\partial
\let\na\nabla 
\let\eps\varepsilon 

\usepackage{color}
\usepackage{caption}

\title[MM\_Propagation userguide]
{Simulation of short pulses propagation in multimode optical fibers}

\author[A. Roget]{Alexandre Roget}
\address[A. Roget]{Inria, Univ. Lille, CNRS, UMR 8524}

\begin{document}

\setlength{\parindent}{0pt}

%%% TITLE AND AMS CLASSIFICATION

\maketitle

\section*{Introduction}
This user guide describes how to use the C++ software \name, developped in the Paradyse
team of Inria by Alexandre Roget.
This code simulates the propagation of short pulses in multimode optical fibers.
Starting from physical parameters of the fiber
(length of the fiber, number of modes of the fiber, dispersion coefficients of the modes,
nonlinear coupling coefficients, Raman factor, {\it etc})
and from an input signal in the fiber, the code computes an approximation of the corresponding
output signal at the end of the fiber.
Numerical discretization parameters (in space and time) can be set by the user as well.
For convenience, the input parameters can be set in a MatLab {\tt .mat} file and
the output is stored in a MatLab {\tt .mat} file as well.

This user guide is organised as follows.
In section \ref{sec:equations}, we describe the physical model solved by \name.
In section \ref{sec:method}, we describe the numerical method used by \name.
In section \ref{sec:setup}, we describe the libraries required to use \name and how to compile the C++ code.
In section \ref{sec:examples}, we provide a set of data to illustrate on an example
how the \name
code works.
In section \ref{sec:comparison}, we compare the performance of \name
with that of other existing software for the same problem.

\section{The physical model}
\label{sec:equations}

The physical model consists in the set of $M\geq 1$ coupled nonlinear Schr\"odinger
equations with nonlocal terms incorporating the Raman effect, of the form

\begin{eqnarray}
  \label{eq:original}
  \lefteqn{\partial_Z \Phi_{p}(Z,T) =
  \ii\beta_0^{(p)} \Phi_p
  -\beta_1^{(p)} \partial_T \Phi_p
  + \ii\frac{\beta_2^{(p)}}{2} \partial_T^2\Phi_p(Z,T)} \\
  & & +\ii\gamma\sum_{k,\ell,m} Q^{(p)}_{k,\ell,m}
  \left(\Phi_k (Z,T) \left[(1-f_R)\Phi_\ell (Z,T)\overline{\Phi_m}(Z,T)
  + f_R\left(h_R\star (\Phi_\ell\overline{\Phi_m})\right)(Z,T)\right]\right) \nonumber
\end{eqnarray}

$\forall p\in\{1,\dots,M\}$ \\

where $\Phi_{p}$ is the field envelope of the $p$th mode, $T$ is the time in a reference frame, $Z$ is the distance along the fiber, $(\beta_0^{(p)})$, $(\beta_1^{(p)})$ and $(\beta_2^{(p)})$ are the dispersion coefficients for the $p$th mode, $\gamma$ is the nonlinearity constant, $f_R\in[0,1]$ is the fractional contribution of the Raman response, the $M^3$ real numbers $(Q^{(p)}_{k,\ell,m})$ are the coupling coefficients for the $p$th mode, and $\tau \mapsto h_R(\tau)$ is the delayed Raman response function that can be written in the form:

\begin{equation}
  \label{eq:raman}
  h_R(t) = \left( \tau_1^{-2} + \tau_2^{-2} \right) \tau_1 {\rm e}^{(-t/\tau_2)} sin(t/\tau_1)
\end{equation}
where $\tau_1$ and $\tau_2$ are the damping time of vibrations. \\

The Lawson-RK method operates with dimensionless values. Thus a conversion is needed. \\

We assume that the pulse has:
\begin{itemize}
  \item peak power $P_0>0$,
  \item pulse width $T_0>0$,
\end{itemize}
and set, as is classical
\begin{equation*}
  L_D = \frac{T_0^2}{\max_p |\beta_2^p|},
  \qquad
  L_{NL} = \frac{1}{\gamma P_0}
  \qquad \text{and} \qquad
  N^2 = \frac{L_D}{L_{NL}} = \frac{T_0^2 \gamma P_0}{\max_p |\beta_2^p|}
  \qquad \text{and} \qquad
  \mathcal Q = \max_{k\ell,m,p} |Q_{k\,ell,m}^{(p)}|.
\end{equation*}
Then, letting
\begin{equation*}
  z=\frac{Z}{L_D},
  \qquad
  t=\frac{T}{T_0},
  \qquad\text{and}\qquad
  \varphi_p (z,t) = \sqrt{|\gamma| L_D} \times \sqrt{\mathcal Q} \times \Phi_p(Z,T),
\end{equation*}
one obtains that equation \eqref{eq:original} is equivalent to
\begin{eqnarray}
  \label{eq:originaldimensionless}
  \lefteqn{\partial_z \varphi_p(z,t) =
  i \frac{T_0^2}{\max_q |\beta_2^{(q)}|} \beta_0^{(p)} \varphi_p(z,t)
  - i \frac{T_0}{\max_q |\beta_2^{(q)}|} \beta_1^{(p)} \partial_t \varphi_p(z,t)
  + i \frac12 \frac{\beta_2^{(p)}}{\max_q |\beta_2^{(q)}|}\partial_t^2 \varphi_p(z,t)} \\
  & & +\ii \text{sign}(\gamma) \sum_{k,\ell,m} \frac{Q^{(p)}_{k,\ell,m}}{\mathcal Q}
  \left(\varphi_k (z,t) \left[(1-f_R)\varphi_\ell (z,t)\overline{\varphi_m}(z,t) 
  + f_RT_0\left(h_R^{T_0}\star(\varphi_\ell\overline{\varphi_m})\right)(z,t)\right]\right) \nonumber
\end{eqnarray}
where the notations are introduced below.

\section{LawsonRK Method}
\label{sec:method}

We assume periodic boundary conditions in time, and work with $k\in\Z$ as the dual Fourier variable
of time. The symbol of the operator $L_p$ defined by
\begin{equation}
  \label{eq:defLp}
  L_p \varphi =
   \ii\beta_0^{(p)} \varphi
  -\beta_1^{(p)} \partial_t \varphi
  +\ii\frac{\beta_2^{(p)}}{2} \partial_t^2\varphi,
\end{equation}
is therefore equal to
\begin{equation}
  \label{eq:symboleLp}
  \mathcal L_p(k) = \ii\beta_0^{(p)}-\ii\beta^{(p)}_1 \alpha(k)-\ii\frac{\beta_2^{(p)}}{2}\alpha^2(k),
\end{equation}
where $-\alpha(k)^2$ is the $k^{\rm th}$ eigenvalue of the Laplace operator with periodic boundary
conditions ({\it e.g.} $\alpha(k)=\frac{2\pi}{T}k$ if $T$ denotes the time period).
Note that $\mathcal L$ has only purely imaginary values; therefore it is anti-symmetric.

Setting for all $p\in\{1,\dots,M\}$,
\begin{equation}
  \label{eq:chgtvar}
  \psi_p(z,t) = {\rm e}^{-zL_p} \varphi_p,
\end{equation}
we obtain the following equivalent equations for \eqref{eq:originaldimensionless}
\begin{equation}
  \label{eq:newpsi}
  \resizebox{0.94\hsize}{!}{$
  \partial_z \psi_p = \ii\gamma {\rm e}^{-z L_p} \sum_{k,\ell,m} Q^{(p)}_{k,\ell,m}
  \left( {\rm e}^{z L_k}\psi_k \right) \left[
    (1-f_R) \left(
      \left( {\rm e}^{z L_\ell} \psi_\ell \right)
      \overline{ \left( {\rm e}^{z L_m} \psi_m \right) } \right)
    + f_R \left(
      h_R \star \left(
        \left( {\rm e}^{z L_\ell} \psi_\ell \right)
        \overline{ \left( {\rm e}^{z L_m} \psi_m \right) } \right)
    \right)
  \right] $}
\end{equation}

Putting the $M$ unknowns in a column-vector $\Psi$ allows to write \eqref{eq:newpsi} under the form
\begin{equation}
  \label{eq:psigroupe}
  \partial_z \Psi(z,t) = N(z,\Psi(z,t)).
\end{equation}

Assume $s\in\N^\star$ is given and an explicit Runge--Kutta method with coefficients
$(a_{i,j})_{1\leq j<i\leq s}$, $(c_{i})_{1\leq i\leq s}$ and $(b_i)_{1\leq i\leq s}$ is given.
Applying it to \eqref{eq:psigroupe} yields to the following explicit computations starting
from $\Psi_n$.
For all $i\in\{1,\dots,s\}$, compute
\begin{equation}
  \label{eq:RKintern}
  \Psi_{n,i}=\Psi_n+h\sum_{j=1}^{i-1} a_{i,j} N(z_n+c_j h,\Psi_{n,j}),
\end{equation}
(with $0$ as a sum when $i=1$), and then compute
\begin{equation}
  \label{eq:RKexplicit}
  \Psi_{n+1} = \Psi_n + h\sum_{i=1}^s b_i N(z_n+c_i h, \Psi_{n,i}).
\end{equation}

We implemented the method above involving the celebrated RK4 method, for which $s=4$ and the order is $4$. Its Butcher tableau is given by
\begin{equation*}
  \renewcommand\arraystretch{1.2}
  \begin{array}
  {c|cccc}
    0 & 0 & 0 & 0 & 0\\
    \frac{1}{2} &  \frac{1}{2} & 0 & 0 & 0\\
    \frac{1}{2} & 0 & \frac12 & 0 & 0\\
    1 & 0 & 0 & 1 & 0  \\
    \hline
    & \frac{1}{6} & \frac{1}{3} & \frac13 & \frac16
  \end{array}.
\end{equation*}

Another class of $6$-order explicit Runge--Kutta methods with $s=7$ stages is given \cite{Butcher64} by
\begin{equation*}
  \begin{array}{c|ccccccc}
    0 & 0 & 0 & 0 & 0 & 0 & 0 & 0 \\
    \mu & \mu & 0 & 0 & 0 & 0 & 0 & 0 \\
    \frac23 & \frac23-\frac{2}{9\mu} & \frac{2}{9\mu} & 0 & 0 & 0 & 0 & 0 \\
    \frac13 & \frac{5}{12}-\frac{1}{9\mu} & \frac{1}{9\mu} & -\frac{1}{12}  & 0 & 0 & 0 & 0\\
    \frac12 & \frac{17}{16}-\frac{3}{8\mu} & \frac{3}{8\mu} & -\frac{3}{16} & \frac{-3}{8} & 0 & 0 & 0\\
    \frac12 & \frac{17}{16}-\frac{3}{8\mu} + \frac{1}{4\lambda} & \frac{3}{8\mu} & - \frac{3}{16} - \frac{3}{4\lambda} & -\frac{3}{8}-\frac{3}{2\lambda} & \frac{2}{\lambda} & 0 & 0\\
    1 & -\frac{27}{44} + \frac{3}{11\mu} & -\frac{3}{11\mu} & \frac{63}{44} & \frac{18}{11} & 4\frac{\lambda-4}{11} & -\frac{4\lambda}{11} & 0 \\
    \hline
    & \frac{1}{120} & 0 & \frac{27}{40} & \frac{27}{40} & \frac{\lambda-8}{15} & -\frac{\lambda}{15} & \frac{11}{120}
  \end{array},
\end{equation*}
for $\lambda,\mu\in\R\setminus\{0\}$.

In the end, after $N\in\N^\star$ "time" steps in the fiber,
one can come back to the original variables by inverting \eqref{eq:chgtvar} and computing for all $p\in\{1,\dots,M\}$,
\begin{equation*}
  \varphi_N^{(p)} = {\rm e}^{z_N L_p} \psi_N^{(p)}.
\end{equation*}

\section{Setup}
\label{sec:setup}

In this section, we describe the libraries required to compile \name,
and how to compile it on your computer.

\subsection{Compiler}

The \name code is compiled with the GNU C++ compiler g++. 
Check g++ compiler version information:     

{\tt \$ g++ --version}

If your system doesn’t have the g++ compiler, you can type the following command in your terminal:

{\tt \$ sudo apt install g++}

\subsection{FFTW}

The numerical one step method described in section \ref{sec:method}
is coupled with a uniform time discretization and the action of dispersive operators
along the fiber is computed {\it via} fast Fourier transforms.
To do so, \name
requires the \href{http://www.fftw.org/fftw2\_doc/fftw\_6.html}{FFTW},
which is a C subroutine library for computing the Discrete Fourier Transform (DFT)
in one or more dimensions, of arbitrary input size, and of both real and complex data. Get the latest stable release of FFTW on the \href{http://www.fftw.org/download.html}{download page}.

\subsection{OpenMP}

The compilation of \name
also requires the application programming interface \href{https://www.openmp.org/}{OpenMP}
that supports multi-platform shared-memory multiprocessing. If OpenMP is not featured in the compiler, you can configure it by running the following command in your terminal:

{\tt \$ sudo apt install libomp-dev}

\subsection{MATIO}

The compilation of \name
requires the MATIO library
that enables to read and write binary MATLAB {\tt .mat} files.
The building procedure of MATIO is described in the \href{https://github.com/tbeu/matio}{github}.

\subsection{MatLab}

For convenience, inputs and outputs for \name
are done using MatLab {\tt .mat} files.
Therefore, a basic version of MatLab can be helpful (no specific package or toolbox is required).
And the {\tt .mat} input files can be generated without MatLab if you prefer.

\subsection{Compilation of MM\_Propagation}
\label{subsec:compilation}
Once all libraries are installed (see the previous subsections), one may write in a terminal:

{\tt \$ git clone} https://github.com/alexandreroget/MM_Propagation \\
{\tt \$ mkdir build} \\
{\tt \$ cd build} \\
{\tt \$ cmake..} \\
{\tt \$ make} \\
{\tt \$ sudo make install} \\

The installation is now complete.

\section{Use MM\_Propagation}
\label{sec:run}

We describe in this section how to set the physical and numerical parameters for the computation
of an approximated output in the fiber by \name  (subsection \ref{subsec:parameters})
and then how to run \name  with these input parameters (subsection \ref{subsec:howtorun}).


\subsection{Set input parameters}
\label{subsec:parameters}

Before doing any propagation simulations, input parameters have to be set.
The \name code gets them from a {\tt .mat} file. Thus open MATLAB and set the following parameters:

\begin{tabular}{|l|c|c|c|c|c|}
  \hline
  \bf{Parameter} & \bf{Notation} & \bf{Unit} & \bf{Variable name} & \bf{Data type} & \bf{Size} \\
  \hline
  Number of modes & $M$ & & n\_modes & integer & 1 \\
  \hline
  Fiber length & $L_0$ & $m$ & fiber\_length & float & 1 \\
  \hline
  Dispersion order & $D$ & & max\_dispersion\_order & integer & 1 \\
  \hline
  Dispersion coefficients & $\beta_i^{(p)}$ & $ps^2/m$ & dispersion\_coefficients & float & $M \times D$ \\
  \hline
  Coupling coefficients & $Q_{k,\ell,m}^{(p)}$ & $m^{-2}$ & coupling\_coefficients & float & $M^4$ \\
  \hline
  Number of points & $N_T$ & & N & integer & 1 \\
  \hline
  Time window & $T_{final}$ & $ps$ & time\_window & float & 1 \\
  \hline
  Pulse width & $T_0$ & $ps$ & pulse\_width & float & 1 \\
  \hline
  Input signal & $\Phi^{(p)}$ & W & initial\_condition & complex & $M \times N_T$ \\
  \hline
  Number of steps & $N_Z$ & $m$ & n\_steps & float & 1 \\
  \hline
  Runge-Kutta order & $s$ & & order\_RK & int & 1 \\
  \hline
  Nonlinearity constant & $\gamma$ & $W^{-1} m$ & nonlinearity\_const & float & 1 \\
  \hline
  Raman proportial & $f_R$ & & raman\_proportial & float & 1 \\
  \hline
  Raman parameters & $\tau_1$ and $\tau_2$ & $ps$ & raman\_parameters & float & 2 \\
  \hline
\end{tabular}

Save input parameters in a {\tt .mat} file with the {\tt save} command.
The \texttt{readMatFile} function of the \name code reads this file and convert MATLAB data. \\

Variables names must be identical to those from the above tabular. Otherwise the function will not be able to read all input parameters.
Arrays sizes must be valid too. For example, if the \texttt{initial\_field} array size is not equal to $M \times N$, then an error will occur and the program will not compute the pulse propagation. \\

The {\tt setParameters.m} file shows an example of how to set input parameters.
You can edit this file and adjust parameters for your simulation, or take it as a template.

\subsection{Run MM\_Propagation}
\label{subsec:howtorun}

Once all input parameters have been set correctly, you can run the \name program. The number of threads \texttt{N\_THREADS} can bet set by typing in a terminal: \\

{\tt \$ export OMP\_NUM\_THREADS = N\_THREADS} \\

Then, type the following command in the terminal: \\

{\tt \$ ./MM\_Propagation input\_parameters.mat output.mat} \\

An error will occur if the input parameters {\tt .mat} file is not specified.
On the other hand, the output filename is optional.
If its filename and its repertory are not specified,
then the simulation results are saved by default in the {\tt output.mat} file.

\section*{Contact}

Technical questions may be directed to Alexandre Roget (alexandre.roget@inria.fr).

\end{document}
